\documentclass[12pt]{article}
\usepackage[utf8]{inputenc}
\usepackage[T1]{fontenc}
\usepackage{hyperref}
\title{Wambaya Language Resources for K-12 Curriculum Development}
\author{Compiled by Claude Sonnet 5.5 (claude-sonnet-5-5) for Andr\'as Kornai}
\date{}
\begin{document}
\maketitle

\section{Introduction}
Wambaya is an endangered Aboriginal language of the Barkly Tablelands region of the Northern Territory, Australia (Mirndi family). Search results gave a 2016 census count of 61 people speaking Wambaya at home. I found no evidence of a school curriculum or published teaching materials. Little exists beyond linguistic documentation.

\section{Primary Educational Resources}
I found no primers, textbooks or school programmes for Wambaya in search results. The Endangered Languages Project entry says Nordlinger's grammar was written in the hope of helping younger Wambaya speakers learn something of their language, but it is a linguistic description and not a teaching course.
\begin{itemize}
\item Endangered Languages Project entry for Wambaya: \url{https://www.endangeredlanguages.com/node/108885}
\end{itemize}

\section{Linguistic and Reference Materials}
\begin{itemize}
\item Rachel Nordlinger, \emph{A Grammar of Wambaya, Northern Territory (Australia)}, Pacific Linguistics, 1998. Record: \url{https://findanexpert.unimelb.edu.au/scholarlywork/1198607-a-grammar-of-wambaya--northern-australia}. The Endangered Languages Project entry says a full-text PDF is hosted on the University of Melbourne Minerva Access repository. I did not see that PDF's URL, so it is not listed in the json.
\item Nordlinger's 1993 MA thesis, \emph{A grammar of Wambaya} (University of Melbourne), includes a short wordlist. I saw no URL for it.
\item ANU list of Central NT languages resources, including papers by Nordlinger: \url{https://www0.anu.edu.au/linguistics/nash/aust/resources.html}
\item Emily Bender, 2008, ``Evaluating a crosslinguistic grammar resource: a case study of Wambaya'' (ACL-08). I saw no URL for it.
\end{itemize}
No dictionary turned up in the searches.

\section{Religious and Community Resources}
General encyclopedic background on the Wambaya people: \url{https://en.wikipedia.org/wiki/Wambaya_people}. I found no scripture translations and no community language organisation in the search results.

\section{Digital and Audio Resources}
General language overview: \url{https://en.wikipedia.org/wiki/Wambaya_language}. I found no audio, video or corpus resources in the search results.

\section{Curriculum Development Context}
Wambaya has few speakers and the only substantial documentation is the Nordlinger grammar. A curriculum would have to be built from that grammar together with new work with community elders. Searches found no ready-made classroom materials.

\end{document}